\section{Theoretical Setup}\label{sec:setup}

\subsection{Problem Setting}

Let $X$ be an arbitrary domain, possibly empty, equipped with its power-set
sigma algebra, and let $H\subseteq\{+1,-1\}^{X}$ be nonempty.  All random
variables and protocol selectors below are measurable.  Since the functions
being integrated are bounded, every displayed expectation is well defined.
For a distribution $D$ on $X$, a target $h\in H$, and a binary predictor
$g:X\to\{+1,-1\}$, define
\begin{equation}\label{eq:loss-definition}
\mathcal L_{D,h}(g)
:=D\{x:g(x)\ne h(x)\}
=\frac{1-\mathbb E_{x\sim D}[h(x)g(x)]}{2}.
\end{equation}
The deterministic dimension complexity is
\begin{equation}\label{eq:dc-definition}
\operatorname{dc}(H):=
\inf\left\{d\in\mathbb N_0:
\begin{array}{l}
\exists\phi:X\to\mathbb R^d,\quad
\forall h\in H\ \exists w_h\in\mathbb R^d,\\[-1mm]
\forall x\in X,\quad h(x)\langle w_h,\phi(x)\rangle>0
\end{array}\right\},
\end{equation}
where $\inf\varnothing:=+\infty$.  Thus one deterministic map must work for
the entire class, and its signs must be strict at every point of the original
domain.  This is neither an average nor a probabilistic,
distribution-dependent, or approximate representation.

Fix
\[
m\in\mathbb N_0,\qquad \tau>0,\qquad
0\le\varepsilon<\frac14,\qquad
\rho:=1-2\varepsilon\in\left(\frac12,1\right],
\]
together with catalog-family constants $B\ge1$ and
$k\in\mathbb N$, $k\ge1$.  A randomized learner $A$ has a random tape
$U$ with fixed law $\nu$ and is deterministic conditional on $U=u$.  At a
nonterminal round $t\le m$, after the full real reply transcript
$v_{<t}\in\mathbb R^{t-1}$, it issues a query
\[
q_{u,v_{<t}}:X\times\{+1,-1\}\longrightarrow[-1,1]
\]
that may depend arbitrarily on $u$ and $v_{<t}$.  It may stop at any depth
$T\le m$, including $T=0$, and then returns a binary predictor.  Every query
has the canonical decomposition
\[
a_q(x):=\frac{q(x,+1)+q(x,-1)}2,qquad
b_q(x):=\frac{q(x,+1)-q(x,-1)}2,
\]
so that $q(x,y)=a_q(x)+yb_q(x)$ and
$|a_q(x)|+|b_q(x)|\le1$.  Neither component is restricted; in particular,
label-independent query components are allowed.

For $D,h,q$, define the population center and full tolerance interval by
\[
\mu_q(D,h):=\mathbb E_{x\sim D}q(x,h(x)),qquad
I_q(D,h):=[\mu_q(D,h)-\tau,\mu_q(D,h)+\tau].
\]
A valid adaptive policy $\pi\in\Pi(D,h)$ is nonanticipating.  After observing
all preceding query-response pairs and the currently issued query $q_t$, it
may choose any real response
\[
\begin{aligned}
v_t&=\pi_t(q_1,v_1,\ldots,q_{t-1},v_{t-1},q_t),\\
v_t&\in I_{q_t}(D,h).
\end{aligned}
\]
The policy may react to every revealed randomized query but not to unrevealed
future learner coins.  A randomized oracle is covered tape by tape after its
own seed is fixed, so the accuracy premise below averages only over learner
coins.  No response alphabet, grid, exact-expectation response, or favorable
policy is imposed.  Write $e(u;D,h,\pi)$ for the complete execution and
$\widehat h_{u;D,h,\pi}$ for its terminal predictor.

\subsection{Assumptions}

\begin{assumption}[Boundary-inclusive source regime]
\label{assump:source-parameter-regime}
The parameters satisfy $m\in\mathbb N_0$, every finite $\tau>0$ is allowed,
and $0\le\varepsilon<1/4$.  In particular, zero-query protocols, exact
learning $(\varepsilon=0)$, and large tolerances are included; no condition
$m\ge1$, $\tau\le1$, or $\varepsilon>0$ is imposed.
\end{assumption}

\begin{assumption}[One randomized adaptive learner]
\label{assump:finite-horizon-randomized-adaptivity}
The same learner $A$, fixed before $D$ and $h$, uses $U\sim\nu$, may choose
each query from the random tape and the entire preceding real-valued reply
transcript, and terminates after at most $m$ queries.  It is not changed after
observing the distribution, target, or oracle policy.
\end{assumption}

\begin{assumption}[Bounded unrestricted SQ access]
\label{assump:bounded-unrestricted-queries}
Every issued query is an arbitrary map
$q:X\times\{+1,-1\}\to[-1,1]$.  No correlational-SQ, label-linearity,
finite-dictionary, nonadaptivity, deterministic-query, finite-encoding, or
response-tree-size restriction is imposed.
\end{assumption}

\begin{assumption}[All continuous tolerance-valid replies]
\label{assump:full-adversarial-tolerance}
For every $D$ and $h$, the learner is evaluated against every
nonanticipating policy in $\Pi(D,h)$, with every response chosen arbitrarily
from the full real interval $I_q(D,h)$.  Exact population responses and
favorable response selection are not assumed.
\end{assumption}

\begin{assumption}[Static terminal-output factorization]
\label{assump:finite-terminal-catalog}
The learner specification exposes one finite catalog
$G=\{g_1,\ldots,g_L\}\subseteq\{+1,-1\}^{X}$ with $L\ge1$, fixed before
$D,h$, the reply policy, all replies, and the sampled learner tape.  On every
complete tolerance-valid execution there is an explicit measurable selector
$J$ such that
\[
\widehat h_{u;D,h,\pi}=g_{J(e(u;D,h,\pi))}
\qquad\text{for every }D,h,\pi,u.
\]
The selector may factor an uncountable continuous transcript tree through
$[L]$; neither a finite response alphabet nor finitely many transcript leaves
is required.  The indexed members $g_1,\ldots,g_L$ are distinct functions.
\end{assumption}

\begin{assumption}[Distribution-target-policy uniform accuracy]
\label{assump:universal-expected-accuracy}
For every distribution $D$ on $X$, every $h\in H$, and every
$\pi\in\Pi(D,h)$,
\[
\mathbb E_{u\sim\nu}
\mathcal L_{D,h}(\widehat h_{u;D,h,\pi})\le\varepsilon.
\]
The displayed expectation is only over the learner tape; distributions,
targets, and tolerance-valid adversarial reply policies are universally
quantified.
\end{assumption}

\begin{assumption}[Boundary-corrected polynomial output budget]
\label{assump:polynomial-catalog-budget}
The catalog size obeys
\[
1\le L\le B\left(1+\frac{m}{\tau^2}\right)^k.
\]
The constants $B\ge1$ and integer $k\ge1$ are fixed for the designated
learner family and are independent of
$X,H,m,\tau,\varepsilon,D,h$, the reply policy and replies, and the learner
coins.  The bound has no hidden constants or suppressed parameter dependence.
\end{assumption}

The catalog may contain functions outside $H$; it constrains terminal outputs
and is not itself a representation hypothesis about $H$.  The proof will
derive, rather than assume, two facts.  First, for every $D,h,\pi$, the output
law on $G$ has probabilities
\[
p_i(D,h,\pi):=\nu\{u:\widehat h_{u;D,h,\pi}=g_i\}
\]
whose mixture correlation is at least $\rho$, and hence some catalog
coordinate has correlation at least $\rho$.  Second, for every finite
$F\subseteq X$, the set
\[
K_{h,F}:=\left\{w\in\Delta_L:
h(x)\sum_{i=1}^Lw_i g_i(x)\ge\rho\ \text{for every }x\in F\right\}
\]
is nonempty, and these sets have the finite-intersection property in the one
fixed simplex $\Delta_L$.  These are proof obligations, not additional
theorem-facing assumptions.

\subsection{Formalized Goal}

The goal is the following exact conditional conclusion.  The one deterministic
catalog map
\[
\phi_G:X\to\mathbb R^L,qquad
\phi_G(x)=(g_1(x),\ldots,g_L(x)),
\]
must satisfy
\[
\forall h\in H\ \exists w_h\in\Delta_L\ \forall x\in X,qquad
h(x)\langle w_h,\phi_G(x)\rangle
\ge\rho=1-2\varepsilon>\frac12.
\]
The map is independent of $D,h$, every valid reply policy and transcript, and
the learner tape.  The weight $w_h$ may depend only on $h$, the fixed catalog,
and the fixed accuracy parameter.  Exact pointwise signs and
\[
\operatorname{dc}(H)\le L
\le B\left(1+\frac{m}{\tau^2}\right)^k
\]
must follow deterministically on every domain.  The conclusion includes
$m=0$, every finite $\tau>0$, $L=1$, $B=1$, $\varepsilon=0$, and the empty
domain.  Its progress type is conditional: it depends on
Assumptions~\ref{assump:finite-terminal-catalog} and
\ref{assump:polynomial-catalog-budget}.  It neither asserts the
boundary-invalid rate $Cm/\tau^2$ nor derives a finite catalog for an arbitrary
unrestricted response tree.  Removing the catalog condition while proving a
boundary-corrected catalog-free linear bound remains open.
